% <КОД>-<N> — <что на рисунке>
% Скопируй файл в исходники/, назови по коду темы (например, П5.2_1.tex) и правь.
\documentclass[tikz,border=3pt]{standalone}
\input{../преамбула}
\begin{document}
\begin{tikzpicture}[pf, scale=0.9]
  % --- 1. Точки: правим только координаты ---
  \coordinate (A) at (0,0);
  \coordinate (B) at (4,0);
  \coordinate (C) at (1.5,2.5);

  % --- 2. Линии по ролям: given / main / find / accent / aux ---
  \draw[given] (A) -- (B) -- (C) -- cycle;

  % --- 3. Отметки: tick1/2/3 на отрезках, ang1 / \pfAngTwo / rightang на углах ---

  % --- 4. Подписи точек ---
  \pfPoint{A}{A}{below left}
  \pfPoint{B}{B}{below right}
  \pfPoint{C}{C}{above}
\end{tikzpicture}
\end{document}
